% Part: first-order-logic
% Chapter: completeness
% Section: downward-ls

\documentclass[../../../include/open-logic-section]{subfiles}

\begin{document}

\olfileid{fol}{com}{dls}
\olsection{लोव्हेनहाइम--स्कोलेम प्रमेय}

एखाद्या उपपत्तीला अनंत प्रतिमान असेल, तर तिला फार तर !!{denumerable}
असलेले प्रतिमानही असते, असे लोव्हेनहाइम--स्कोलेम प्रमेय सांगते. या
तथ्याची तात्काळ निष्पत्ती अशी की एखाद्या !!a{structure} चा आकार
!!{nonenumerable} आहे हे प्रथम-कोटी तर्कशास्त्र व्यक्त करू शकत नाही: सर्व
!!{nonenumerable} !!{structure}s मध्ये पूर्त होणारे कोणतेही !!{sentence}
किंवा !!{sentence}s चा कोणताही संच हा एखाद्या !!{enumerable} रचनेतही
पूर्त होतो.

\begin{thm}
\ollabel{thm:downward-ls} $\Gamma$ सुसंगत असेल, तर त्याला
!!a{enumerable} प्रतिमान असते; म्हणजे ज्याचा प्रांत सांत किंवा
!!{denumerable} आहे अशा !!a{structure} मध्ये तो पूर्ततायोग्य असतो.
\end{thm}

\begin{proof}
$\Gamma$ सुसंगत असेल, तर संपूर्णता प्रमेयाच्या सिद्धतेतून मिळणाऱ्या
!!{structure}~$\Struct M$ चा प्रांत~$\Domain{M}$ हा भाषा~$\Lang L$ मधील
पदांच्या संचापेक्षा मोठा नसतो. म्हणून $\Struct M$ फार तर
!!{denumerable} आहे.
\end{proof}

\begin{thm}
\ollabel{noidentity-ls} $\Gamma$ हा एकरूपतेविना प्रथम-कोटी तर्कशास्त्राच्या
भाषेतील !!{sentence}s चा सुसंगत संच असेल, तर त्याला !!a{denumerable}
प्रतिमान असते; म्हणजे ज्याचा प्रांत अनंत आणि !!{enumerable} आहे अशा
!!a{structure} मध्ये तो पूर्ततायोग्य असतो.
\end{thm}

\begin{proof}
$\Gamma$ सुसंगत असेल आणि त्यात एकरूपता येणारी वाक्ये नसतील, तर संपूर्णता
प्रमेयाच्या सिद्धतेतून मिळणाऱ्या !!{structure}~$\Struct M$ चा
प्रांत~$\Domain{M}$ हा भाषा~$\Lang L'$ मधील पदांच्या संचाशी एकरूप असतो.
म्हणून $\Trm[L']$ हे गणनीय अनंत असल्यामुळे $\Struct{M}$ देखील
!!{denumerable} आहे.
\end{proof}

\begin{ex}[स्कोलेमची विरोधापत्ती]
त्सेर्मेलो--फ्रेंकेल संचसिद्धान्त~$\Log{ZFC}$ ही एक अतिशय समर्थ चौकट
आहे; संचांच्या आकारांविषयीच्या तथ्यांसह जवळजवळ सर्व गणिती विधाने तिच्यात
व्यक्त करता येतात. म्हणून, उदाहरणार्थ, वास्तव संख्यांचा संच~$\Real$ हा
!!{nonenumerable} आहे हे~$\Log{ZFC}$ सिद्ध करू शकते; कोणत्याही संचाचा
घातसंच त्या संचापेक्षा मोठा असतो हे कँटर यांचे प्रमेयही ती सिद्ध करू शकते;
आणि असेच पुढे. $\Log{ZFC}$ सुसंगत असेल, तर तिची सर्व प्रतिमाने अनंत
असतात. शिवाय, त्या सर्व प्रतिमानांत असे !!{element}s असतात की उपपत्तीच्या
मते ते !!{nonenumerable} असतात; उदाहरणार्थ, नैसर्गिक संख्यांचा घातसंच
अस्तित्वात आहे हे~$\Log{ZFC}$ चे प्रमेय सत्य करणारा घटक. लोव्हेनहाइम--स्कोलेम
प्रमेयानुसार $\Log{ZFC}$ ला !!{enumerable} प्रतिमानेही असतात---ज्यांत
``!!{nonenumerable}'' संच असतात, पण जी स्वतः !!{enumerable} असतात.
\end{ex}

\end{document}
